\documentclass{article}
\usepackage[T1]{fontenc}
\usepackage{lmodern}
\usepackage{graphicx}
\usepackage{amsmath,amssymb}
\usepackage[a4paper,margin=2cm]{geometry}
\usepackage[absolute,overlay]{textpos}
\setlength{\TPHorizModule}{1mm}
\setlength{\TPVertModule}{1mm}
\usepackage[indonesian]{babel}
\usepackage{hyperref}
\setlength{\emergencystretch}{2em}
% unit: Halaman 34

\begin{document}

% === Halaman 34 (python/native) ===

• Untuk membuat enkripsi menjadi lebih kompleks, gunakan kata kunci sepanjang 

n dengan huruf-huruf berbeda. Urutan huruf di dalam kata kunci menentukan 

urutan pembacaan secara vertikal.

• Contoh: kata kunci = TOMBAK    (urutan huruf sesuai alfabet: 6  5  4  2  1  3)   

 Plainteks: sistem dan teknologi informasi itb

   Enkripsi:   

             123456   

             TOMBAK


sistem


dantek


nologi


inform


asiitb

34

Cipherteks: (baca secara vertikal sesuai urutan 

huruf di dalam kata kunci → 5, 4, 6, 3, 2, 1)

EEGRTTTOOIMKIMBSNLFIIAONSSDNIA

\end{document}
